% Additional manuscript-ready prose. Integrate selectively with other coauthors.

\subsection{Why economic capability and conduct belong in one evaluation}

Persistent economic tasks make planning consequences observable over more than one response. An agent can acquire assets, preserve cash, borrow against property and recover from debt. Competing agents add a second source of consequences. A payment can finance a rival, a rejected offer can preserve a blocking position, and a bankruptcy can transfer productive assets. Vending-Bench motivates sustained economic operation as an agent-evaluation problem, while subsequent multiplayer work studies negotiation and conduct within related economies~\cite{backlund2025vendingbench,li2026arena,wu2026agentmonopoly}. These settings allow capability and social behavior to be examined along the same trajectory without requiring that financial success define acceptable conduct.

MonopolyBench makes the relationship concrete. The environment provides an economic objective, explicit legal choices and records linking decisions to state changes. Its current instructions permit public deception. A financially misleading public bargain can therefore be effective competitive behavior without violating a truthfulness requirement. Other questions remain open to evaluation: whether the proposition is correct, what information the recipient receives, whether the agent keeps a stated commitment, and whether its applied choice agrees with its accompanying account. A statement's factual status, contractual status and effect on a counterparty are distinct empirical objects. Reward and normative conduct also require separate definitions in broader safety evaluations~\cite{pan2023machiavelli,ren2025mask}.

\subsection{Economic consequences do not identify social influence}

The G13 Pennsylvania negotiation shows how economic reconstruction can strengthen a social account while limiting a causal reading. Before the disputed pitch, the buyer at d251 already offers \$130 and records a reason to preserve the seller's solvency because the seller may later pay Short Line rent to the buyer. At d252 the seller presents \$130 as insufficient for its \$200 bill, although \$35 cash, the \$130 payment and a \$130 Ventnor mortgage cover the debt. Its applied private note acknowledges this route, and its returned account discusses withholding it. The buyer's delivered prompt at d253 includes the pitch, and the buyer accepts \$200. The settlement transfers the deed and cash, with \$16 paid to the bank in mortgage-transfer interest.

The sequence supports a delivered misleading omission, a recorded omission rationale and a subsequent higher-price settlement. The buyer's earlier offer and solvency rationale also supply an alternative explanation for acceptance. The extra \$70 cannot be attributed to persuasion from this single sequence. Nor does recorded arithmetic correction establish complete resistance to price framing or other influence. G19 instead pairs explicit premise correction with acceptance, while G11 contains a transient echo of a false premise followed by checking and rejection. Those contrasts describe different recipient responses to economic claims. They do not estimate a general susceptibility rate.

Promises require an equally specific object. In G18, a standing \$250 offer remains in both parties' delivered memories, is explicitly invoked, and is then withdrawn after intervening rent, sales and mortgages. The records establish exposure, invocation and withdrawal in a changed state. They do not establish an irrevocable commitment or a costly action undertaken in reliance. The relevant question is what assurance the parties actually exchanged, with what conditions and release authority. Work on commitments provides useful definitions for designing that question, rather than retrospectively converting every bargaining offer into a promise~\cite{empty_commitments_v2,commitment_standing}.

\subsection{What a capability--alignment comparison would require}

The present observations motivate a controlled comparison, but do not provide its capability axis. Win counts depend on response handling, opportunities, opponents and random paths; they cannot order innate capability and then reveal how alignment changes along that order. Literature on agent conduct already shows that the apparent relationship depends on the task and evaluated norm~\cite{li2026arena,shi2026emergentcollusion,ren2024safetywashing}. Our contribution is to identify measurement conditions that matter before constructing such a comparison.

A prospective study could specify truthful financial disclosure, protected information or a transaction-authorization rule independently of the profit objective. It would then create situations with genuine opportunities both to comply and to violate. An independently measured capability axis, or controlled tiers within a model lineage, would replace the current roster's realized win ordering. Exact state restoration could hold assets, debt and available actions fixed while prompts, budgets and response conditions are controlled. Fresh repeated policy continuations would be required for downstream effects, with held-out situations used to evaluate both task performance and norm adherence. A legal alternative at one historical prefix is a feasibility check within this design, rather than an observed policy under the alternative treatment.

Randomization is not necessary for every descriptive statement of adherence to a specified norm. An observable unauthorized transfer can violate an explicit authorization rule in one trace. Controlled opportunities and repeated continuations become necessary for causal comparisons and generalization. Likewise, private notes and provider explanations can document a stated reason but do not uniquely identify the computation behind a choice~\cite{turpin2023unfaithful,lanham2023faithfulness}. Received channel disagreements warrant scrutiny, but a monitor would need specified inputs, a prediction rule and held-out adjudicated outcomes before these disagreements could become an accuracy measure. Task utility, individual norm violations and their consequences should remain separately reportable rather than being silently combined into a single profit-adjusted alignment score.
